\section{Security}\label{sec:agg-creds-security}

%%%% proof 1 - the generic construction is secure (Theorem 1)

\subsection{Proof of \Cref{thm:agg-creds-construction-security}: security of the construction}\label{sec:agg-creds-proof-1}

We define a sequence of indistinguishable games
that start with Game 0 representing the real world, and end with Game 8
representing the ideal world.

The games use the four hypotheses of
\Cref{thm:agg-creds-construction-security}:

\begin{enumerate}

  \item[(P1)] The registration authority's signature scheme $\DS$ is
existentially unforgeable (EUF-CMA).

  \item[(P2)] The issuers' tag-based aggregate signature scheme $\ATS$ is
existentially unforgeable under tag-based aggregation.

  \item[(P3)] The proof system for $\algoproof$ is zero knowledge and
simulation extractable.

  \item[(P4)] The function $f$ relating users' keys is one way, and the proofs
of knowledge of users' and issuers' secret keys are zero-knowledge proofs of
knowledge.

\end{enumerate}

\begin{proof}

  \paragraph{Game 0.} This is the execution of $\prot$.

  \paragraph{Game 1.} As Game 0, except $\simu$ executes $\prot$ on
behalf of the honest parties: $\simu$ is provided with the inputs (including
secrets) of the honest parties, and interacts with corrupt parties, through the
adversary, in the real world, following $\prot$'s specifications, and returns the
corresponding output. At setup, $\simu$ extracts the secret key of every corrupt
issuer from the proof of knowledge published with its public key. Corruption is
static, so this extraction happens once, before any other interaction. This
change is purely syntactic, so Game 1 is indistinguishable from Game 0 by
construction.

  \paragraph{Game 2.} As Game 1, except we introduce dummy functionality
$\idfu{\AAC}'$ that forwards the inputs from the honest parties to $\simu$ and
returns the outputs. This modification is syntactic, so Game 2 is
indistinguishable from Game 1.

  \paragraph{Game 3.} As Game 2, except $\idfu{\AAC}'$ handles register calls
from honest users following the description of ideal functionality
$\idfu{\AAC}$. $\simu$ first assigns each honest user a key pair $(\sk, \pk)$.
Now, when an honest user $\user$ calls $\idfu{\AAC}'$ to register, $\simu$
learns the user's identity. Given this information, $\simu$ retrieves the
corresponding key pair $(\sk,\pk)$ and contacts the registration authority to
register $\user$. $\simu$ receives a unique tag $\Tag$ and a signature $\psi
\sample \DSsign(\rask, (\pk, \Tag))$ and stores entry $\langle \user, \Tag,
\pk, \psi\rangle $.

  Since the registration authority and $\idfu{\AAC}'$ both register each user
exactly once and reject duplicates, they accept or reject the same requests, so
this game is indistinguishable from Game 2.

  \paragraph{Game 4.} As Game 3, except $\idfu{\AAC}'$ also handles credential
issuance calls from honest users following the description of ideal
functionality $\idfu{\AAC}$.  When  $\idfu{\AAC}'$ receives issuance request
$(\AACissue, \att{i}, \atv{i})$ from honest user $\user$, $\simu$ learns the
triple $(\user, \att{i}, \atv{i})$.  To simulate the corresponding credential
issuance for $\user$ in the real world, $\simu$ retrieves $(\sk, \pk)$ and
$(\Tag, \psi)$.  If the pair $(\Tag, \psi)$ for $\user$ does not exist, meaning
$\user$ did not call register before, then $\simu$ aborts. Otherwise, it
honestly executes $\prot$ with issuer $\issuer{i}$.  There are two cases to
consider.

  \begin{description}

    \item[$\issuer{i}$ is honest.] As per $\prot$'s specification, $\simu$ receives
signature $\signature_i \sample \ATSsign(\isk{i}, \Tag, \atv{i})$ and
records entry $\langle \user, \Tag, \att{i}, \atv{i}, \signature_i \rangle$, and
$\idfu{\AAC}'$ likewise adds $\atv{i}$ to $\creds[\user][i]$.

    \item[$\issuer{i}$ is corrupt.] If the execution of $\prot$ terminates with
$\simu$ receiving a signature $\signature_i$ with $\ATSverify(\ipk{i}, \Tag,
\atv{i}, \signature_i) = 1$, then $\simu$ records entry $\langle \user, \Tag,
\att{i}, \atv{i}, \signature_i \rangle$, and sends accept.

  \end{description}

  In both cases, $\idfu{\AAC}'$ follows $\simu$. An honest issuer accepts every
request from a registered user, matching the first branch, and for a dishonest
issuer, $\idfu{\AAC}'$ accepts only when $\simu$ does, which is when the issuer
produced a valid single $\ATS$ signature.  So $\idfu{\AAC}'$'s output matches
$\prot$'s in both cases, and this game is indistinguishable from Game 3 as long as
$\simu$ does not abort.  Since $\user$ is honest, they always register
before requesting credentials, so the abort event never occurs.

  \paragraph{Game 5.} As Game 4, except $\idfu{\AAC}'$ executes aggregate
requests from honest users following the description of $\idfu{\AAC}$. Since
credential aggregation is a local operation where adversary $\adv$ observes
nothing, $\simu$ need not do anything, so Game 5 is indistinguishable from Game
4.

  \paragraph{Game 6.} As Game 5, except $\idfu{\AAC}'$ takes charge of presentation
requests from honest users, following the specifications of $\idfu{\AAC}$.
Given the message $(\AACpresentend, \vec{\atv{}}, \msg, \algoproof)$ from
$\idfu{\AAC}'$ for $\vec{\atv{}} = (\atv{1}, \ldots, \atv{n})$, $\simu$ defines
the statement $\inst_1 = (\rapk,\allowbreak \ISSUERS,\allowbreak \{\ipk{i}, \atv{i}\}_{i \in \ISSUERS})$ such that $\forall i \in
\ISSUERS: \atv{i} \neq \varnothing$, and, using the simulation trapdoor $\td$
of the proof system, outputs $\algoproof^\star
\sample \NIZKsimprove(\crs, \td, \inst_1, \msg)$ in the real world, which gives
$\NIZKverify(\crs, \inst_1, \algoproof^\star, \msg) = 1$. It records
$\langle \vec{\atv{}}, \msg, \algoproof, \algoproof^\star \rangle$. By (P3), this
simulated proof is indistinguishable from a real one, so Game 6 is
indistinguishable from Game 5.

  \paragraph{Game 7.} As Game 6, except $\idfu{\AAC}'$ processes verification
requests from honest verifiers following the description of $\idfu{\AAC}$, and
$\simu$ simulates the honest verifiers in the real world.  On receiving a
presentation in the real world, $\simu$ translates the presentation into calls to
$\idfu{\AAC}'$, and instead of returning the output of $\NIZKverify$, it returns
the verification result received from $\idfu{\AAC}'$.

  \begin{description}

    \item[Register.] For each registration request $\pk$ that a corrupt user
sends, in the real world, to the honest registration authority, $\simu$ (acting
as the registration authority) executes the registration protocol. If it
succeeds, $\simu$ records entry $\langle\user, \Tag, \pk, \psi
\rangle$, and then (acting as a user) transmits the corresponding registration
request $register$ to $\idfu{\AAC}'$ which processes it in a similar fashion to
$\idfu{\AAC}$.

      The simulated registration authority and $\idfu{\AAC}'$ accept or reject
the same requests by construction, so their outputs are indistinguishable.

    \item[Issue.] For each issuance request $(\Tag, \pk, \psi, \atv{i})$ that a
corrupt user $\user$ sends, in the real world, to an honest issuer $\issuer{i}$
(simulated by $\simu$), $\simu$ (acting as $\user$) sends $(\AACissue, \att{i},
\atv{i})$ to $\idfu{\AAC}'$, and stores entry $\langle \user, \Tag, \att{i},
\atv{i}, \signature_i\rangle$ on success.  $\idfu{\AAC}'$ issues the credential
exactly when $\user$ is registered, and the simulated $\issuer{i}$ returns $\signature_i \sample
\ATSsign(\isk{i}, \Tag, \atv{i})$ under the same conditions, so their outputs
are indistinguishable.

    \item[Verify.] For each presentation $(\vec{\atv{}}, \msg, \algoproof^\star)$,
with $\vec{\atv{}} = (\atv{1}, \ldots, \atv{n})$, that a user $\user$ presents,
in the real world, to an honest verifier $V$, $\simu$ (simulating $V$) calls
$b \leftarrow \NIZKverify(\crs, \inst_1, \algoproof^\star, \msg)$.

      If $b = 0$, $\simu$ sends $(\AACverify, \vec{\atv{}}, \msg,
\algoproof^\star)$ to $\idfu{\AAC}'$ and outputs its answer. The strings of
$\idfu{\AAC}'$ are random, so except with negligible probability the entry of
$\proven$ for $\algoproof^\star$ is $\bot$. Either there is a dishonest user
$\user'$ such that, for every $i$ with $\atv{i} \neq \varnothing$, either
$\atv{i} \in \creds[\user'][i]$ or issuer $\issuer{i}$ is dishonest. Then
$\idfu{\AAC}'$ queries $\simu$, which answers $0$. Or there is no such user, and
$\idfu{\AAC}'$ records the presentation as invalid and returns $0$. In both
cases the output matches that of $\prot$.

      If $b = 1$ and $\simu$ recorded $\langle \vec{\atv{}}, \msg, \algoproof,
\star \rangle$ for an honest user's presentation, then $\simu$ sends
$(\AACverify, \vec{\atv{}}, \msg, \algoproof)$ to $\idfu{\AAC}'$, which
returns $1$, as $V$ does in the real world.

      Otherwise, if $b = 1$, $\simu$ simulated no proof for this statement and
label. It sends $(\AACverify, \vec{\atv{}}, \msg, \algoproof^\star)$ to
$\idfu{\AAC}'$, answers $1$ if $\idfu{\AAC}'$ queries it, and outputs the
answer of $\idfu{\AAC}'$.

  \end{description}

  $\simu$ never extracts, so it runs straight-line. In the last case,
$\idfu{\AAC}'$ queries $\simu$ if there is a dishonest user $\user'$ such that,
for every $i$ with $\atv{i} \neq \varnothing$, either $\atv{i} \in
\creds[\user'][i]$ or issuer $\issuer{i}$ is dishonest. Then $\simu$ answers $1$,
as $V$ does in the real world. If there is no such user, $\idfu{\AAC}'$ returns
$0$, and we call this event $E$. If the same presentation was verified before,
$\idfu{\AAC}'$ returns the earlier output, which is $1$ unless $E$ occurred
then. So this game is indistinguishable from Game 6 as long as $E$ does not
occur.

  We now show that $E$ occurs with negligible probability. For each of the
events $a_1$, $a_2$ and $a_3$ below, a reduction runs Game 7 until the first
verification at which $E$ occurs, and outputs that presentation. It verifies,
and $\simu$ simulated no proof for its statement and label, so by (P3) the
reduction extracts a witness $\witness_1 = (\Tag, \sk, \pk, \psi, \signature)$,
except with negligible probability. This is the only place where the proof
extracts, and each reduction extracts once. Since no dishonest user covers
$\vec{\atv{}}$, one of the following events occurs:

  \begin{description}

    \item[$a_1$.] There is no recorded entry $\langle \star, \Tag, \pk, \star
\rangle$, where $\star$ matches any value. Then $\user$ has knowledge of a
valid $\psi$ such that $1 \leftarrow \DSverify(\rapk, (\pk, \Tag), \psi)$ with
no corresponding entry. The registration authority keeps $\rask$ secret, and
the reduction answers every registration with its signing oracle, so this
contradicts (P1).

    \item[$a_2$.] The recorded entry $\langle \user', \Tag, \pk, \star
\rangle$ belongs to a dishonest user $\user'$. Since $\user'$ does not cover
$\vec{\atv{}}$, there exists $j \in [1, n]$ such that $\atv{j} \neq
\varnothing$, issuer $\issuer{j}$ is honest, and $\atv{j} \notin
\creds[\user'][j]$. Each tag belongs to one user, so there is no recorded entry
$\langle \star, \Tag, \att{j}, \atv{j}, \star \rangle$. The extracted
$\signature$ satisfies the following:
    %
      \begin{equation*}
    %
	\ATSverifyagg((\ipk{i})_{i\in \ISSUERS},\Tag,(\atv{i})_{i\in\ISSUERS},
\signature) = 1.
    %
      \end{equation*}
    %
      $\ISSUERS$ here identifies the indices with $\atv{i} \neq \varnothing$ and
$j \in \ISSUERS$. The honest issuer $\issuer{j}$ never issued a signature
$\signature_j$ for pair $(\Tag, \atv{j})$, contradicting (P2). The reduction to
(P2) sets the key of each honest issuer to one of the challenge keys of
\Cref{exp:ats-euf}, answers its issuance through $\signo$, and corrupts the
other challenge keys. For each corrupt issuer, it outputs the published key and
the secret key extracted in Game 1. The extracted $\signature$ then verifies
under these keys, and $\signo$ was never queried on $(\Tag, \atv{j})$ under the
key of $\issuer{j}$.

    \item[$a_3$.] The recorded entry $\langle \user', \Tag, \pk, \star
\rangle$ belongs to an honest user $\user'$. The extracted $\sk$ is then a
preimage of $\pk$ under $f$. Honest users use $\sk$ only in their proofs of
knowledge at registration and issuance, and their presentations are simulated
since Game 6. So a reduction that simulates these proofs of knowledge, which
are zero knowledge by (P4), inverts $f$ on $\pk$, contradicting its one-wayness.

  \end{description}

  \paragraph{Game 8.} This is the ideal world. As Game 7, except $\idfu{\AAC}'$
is now identical to ideal functionality $\idfu{\AAC}$, so Game 8 is
indistinguishable from Game 7. This concludes the proof
of~\Cref{thm:agg-creds-construction-security}.

\end{proof}

%%%% proof 2 - tag-based agg sigs are secure (Theorem 2)

\subsection{Proof of \Cref{thm:agg-creds-thm-1}: unforgeability of $\ATS$}\label{sec:agg-creds-proof-2}

Let $\pk_1, \ldots, \pk_n$ be the challenge public keys of
\Cref{exp:ats-euf}, and suppose an adversary $\adv$ wins by outputting an
aggregate signature $\signature^* = (\signatureone^*, \signaturetwo^*)$ on
messages $(\msg_i^*)_{i \in \ISSUERS}$ under a tag $\Tag^*$, with keys
$(\pk'_i)_{i \in \ISSUERS}$, so that the following holds:
%
\begin{equation*}
  \pairing(\signaturetwo^*,\ggen) = \pairing(\hashtwo(\Tag^*), \signatureone^*)\cdot
\prod_{i \in \ISSUERS} \pairing(\hashone(\msg_i^*), \pk'_i).
\end{equation*}
%
Some $j \in \ISSUERS$ has an uncorrupted key $\pk_j$, and the signing oracle
never signed $(\pk_j, \Tag^*, \msg_j^*)$. The adversary may have asked for
signatures under $\pk_j$ on $\msg_j^*$ with other tags, and on $\Tag^*$ with
other messages.

\begin{proof}

  We build a reduction $\advB$ that breaks BCDH. It receives a challenge
defined as follows:
%
\begin{equation*}
  (\gen, X, Y, Z, \ggen, \widetilde{X}, \widetilde{Y}, \widetilde{Z}) = (\gen,
\gen^\gamma, \gen^\omega, \gen^\chi, \ggen, \ggen^\gamma, \ggen^\omega,
\ggen^\chi),
\end{equation*}
%
and must output $\pairing(\gen, \ggen)^{\gamma\omega\chi}$.

  \paragraph{Keys.} Without loss of generality, $\adv$ corrupts every key but
one. $\advB$ guesses $j \sample [1,n]$ and sets $\pk_j = \widetilde{X}$. For every
$i \neq j$ it samples $\xi_i \sample \Zp^*$ and sets $\sk_i = \xi_i$ and $\pk_i =
\ggen^{\xi_i}$. It answers a corruption of $\pk_i$ with $\xi_i$, and stops if
$\adv$ corrupts $\pk_j$. For every corrupted key in the forgery, $\adv$ also
outputs its secret key, which we write $\xi'_i$.

  \paragraph{Random oracles.} Let $q$ bound the number of signing queries under
$\pk_j$. $\advB$ keeps lists of the queries to $\hashone$
and $\hashtwo$, with the exponents it used. On a new query $\msg$ to $\hashone$,
it samples $\rho \sample \Zp^*$ and answers $Y^\rho$ with probability $1/(q+1)$, and
$\gen^\rho$ otherwise. We call the message embedded in the first case. On a new
query $\Tag$ to $\hashtwo$, it samples $\lambda \sample \Zp^*$ and answers
$\gen^\lambda$ with probability $1/(q+1)$, and $Y^\lambda$ otherwise. We call the tag plain
in the first case. Every answer is uniform in $\group{1} \setminus \{1\}$.

  \paragraph{Signing.} $\advB$ answers a signing query $(\pk_i, \Tag, \msg)$
with $i \neq j$ by running $\ATSsign(\xi_i, \Tag, \msg)$. For a query under
$\pk_j$, it samples $\beta \sample \Zp$ and proceeds as follows.
  \begin{itemize}

    \item If $\msg$ is not embedded, so that $\hashone(\msg) = \gen^\rho$, it
returns the following signature:
%
\begin{equation*}
  (\signatureone, \signaturetwo) = (\ggen^\beta, \ X^\rho\,\hashtwo(\Tag)^\beta).
\end{equation*}
%
This is $\ATSsign(\gamma, \Tag, \msg)$ with randomness $\beta$, since $X^\rho =
\hashone(\msg)^\gamma$.

    \item If $\msg$ is embedded and $\Tag$ is not plain, so that $\hashone(\msg) =
Y^\rho$ and $\hashtwo(\Tag) = Y^\lambda$, it returns the following signature:
%
\begin{equation*}
  (\signatureone, \signaturetwo) = (\ggen^\beta\,\widetilde{X}^{-\rho/\lambda}, \
Y^{\lambda\beta}).
\end{equation*}
%
This is $\ATSsign(\gamma, \Tag, \msg)$ with randomness $\upsilon = \beta -
\rho\gamma/\lambda$, which is uniform, since the following holds:
%
\begin{equation*}
  \hashone(\msg)^\gamma \hashtwo(\Tag)^{\upsilon} = Y^{\rho\gamma + \lambda\beta -
\rho\gamma} = Y^{\lambda\beta}.
\end{equation*}

    \item If $\msg$ is embedded and $\Tag$ is plain, it stops.

  \end{itemize}

  \paragraph{Extraction.} If $\adv$ never queried $\hashone(\msg^*_j)$ or
$\hashtwo(\Tag^*)$, $\advB$ first queries them itself, so that each is embedded
or plain as above. $\advB$ stops unless $\msg^*_j$ is embedded and
$\Tag^*$ is plain, so that $\hashone(\msg^*_j) = Y^{\rho^*}$ and
$\hashtwo(\Tag^*) = \gen^{\lambda^*}$. It removes the contributions of the other
keys as follows:
%
\begin{equation*}
  d = \signaturetwo^* / \prod_{i \in \ISSUERS, i \neq j}\hashone(\msg^*_i)^{\xi_i},
\end{equation*}
%
with $\xi_i$ replaced by $\xi'_i$ for the corrupted keys. The verification
equation then becomes the following:
%
\begin{equation*}
  \pairing(d, \ggen) = \pairing(\gen^{\lambda^*}, \signatureone^*) \cdot
\pairing(Y^{\rho^*}, \widetilde{X}).
\end{equation*}
%
Pairing both sides with $\widetilde{Z}$ instead of $\ggen$, $\advB$ outputs the
following:
%
\begin{equation*}
  \left(\pairing(d, \widetilde{Z}) / \pairing(Z^{\lambda^*},
\signatureone^*)\right)^{1/\rho^*} = \pairing(Y, \widetilde{X})^{\chi} =
\pairing(\gen, \ggen)^{\gamma\omega\chi}.
\end{equation*}

  \paragraph{Success probability.} The simulation is perfect up to the
difference between uniform elements of $\group{1}$ and of $\group{1} \setminus
\{1\}$, at most $q_H/p$ for $q_H$ hash queries.
$\advB$ does not stop when all of the following
hold:
  \begin{itemize}
    \item the guess of $j$ is right, with probability $1/n$;
    \item $\msg^*_j$ is embedded, with probability $\frac{1}{q+1}$, and none of
the at most $q$ other messages signed under $\pk_j$ is, each independently with
probability $1-\frac{1}{q+1}$, so together with probability at least
$\frac{1}{q+1}\big(1-\frac{1}{q+1}\big)^q$;
    \item $\Tag^*$ is plain, with probability $\frac{1}{q+1}$, and none of the at
most $q$ other tags signed under $\pk_j$ is, each independently with probability
$1-\frac{1}{q+1}$, so together with probability at least
$\frac{1}{q+1}\big(1-\frac{1}{q+1}\big)^q$.
  \end{itemize}
The probability $\frac{1}{q+1}$ maximises $x(1-x)^q$ over $x \in [0, 1]$, which
is why $\advB$ embeds with it.
Every signing query under $\pk_j$ then falls into one of the first two cases. A
query on a message other than $\msg^*_j$ is not embedded, and a query on
$\msg^*_j$ has a tag other than $\Tag^*$, since $(\pk_j, \Tag^*, \msg^*_j)$ was
never queried, so its tag is not plain. The events are independent of the view
of $\adv$, so $\advB$ succeeds with probability at least $\frac{1}{n}$ times
the square of this bound, times that of $\adv$, which gives the following:
%
\begin{equation*}
  \Adv^{\text{BCDH}}_{\advB}(\secparam) \geq \frac{1}{n}
\left(\frac{1}{q+1}\Big(1-\frac{1}{q+1}\Big)^q\right)^2
\left(\Adv^{\mathsf{EUF\text{-}TBA}}_{\ATS,\adv}(\secparam) - q_H/p\right).
\end{equation*}
%
So $\ATS$ is existentially unforgeable under tag-based aggregation.

\end{proof}

%%%% proof of the lemma - extraction and simulation of presentations

\subsection{Proof of \Cref{lem:agg-creds-sigma}: extraction and simulation of presentations}\label{sec:agg-creds-proof-lemma}

\begin{proof}

    Recall the instance $\inst_2$, witness $\witness_2$ and relation $\rel_2$
of \eqref{eq:agg-creds-proof-original}, and $\inst_3$ of
\eqref{eq:agg-creds-proof-rel}. A presentation is $\algoproof = (\eta,
\inst_3)$.

  \paragraph{Extraction.} By special soundness of the Sigma protocol, which
Fiat-Shamir preserves in the random oracle model, if a user is able to produce
a valid proof $\eta$ for relation $\rel_3$, then it
is possible to extract the corresponding witness $\witness_3 = (\delta, \sk)$
satisfying the following:
%
  \begin{equation*}
%
    \bar{\gen} = \gen^\delta, \quad \bar{v} = v^\delta, \quad H_i =
\hashone(\atv{i})^\delta, \quad \pk' = \bar{\gen}^\sk.
%
  \end{equation*}
%
  Since $\bar{\gen} \neq 1$, we have $\delta \neq 0$, and can write the
following:
%
  \begin{equation*}
%
    \gen = \bar{\gen}^{1/\delta}, \quad v = \bar{v}^{1/\delta}, \quad
\pk'^{1/\delta} = \bar{\gen}^{\sk/\delta} = \gen^\sk.
%
  \end{equation*}
%
  A valid presentation satisfies the blinded
$\ATSverifyagg$ and $\SPSverify$ equations:
%
  \begin{gather*}
%
    \pairing(\bar{\signaturetwo}, \ggen) =
\prod_{i\in\ISSUERS}\pairing(H_i,\ipk{i}) \cdot
\pairing(\bar{\tid},\signatureone') \ \wedge \\ \pairing(\bar{\gen},\ggen) =
\pairing(\bar{z},\widetilde{r}')\cdot \pairing(\pk',\widetilde{u}_1) \cdot
\pairing(\bar{\tid},\widetilde{u}_2) \ \wedge \\
\pairing(\bar{v},\widetilde{r}') = \pairing(\bar{\gen},\widetilde{s}').
%
  \end{gather*}
%
  Substituting $\bar{\gen}^{1/\delta} = \gen$ and $\bar{v}^{1/\delta} = v$, and
raising all three equations to the power $1/\delta$, we deduce the following:
%
  \begin{gather*}
%
    \pairing(\bar{v}^{1/\delta}, \widetilde{r}') = \pairing(v, \widetilde{r}') =
\pairing(\bar{\gen}^{1/\delta}, \widetilde{s}') = \pairing(\gen, \widetilde{s}')
\ \wedge \\
  %
    \pairing(\bar{z}^{1/\delta}, \widetilde{r}') \cdot \pairing( \pk'^{1/\delta},
\widetilde{u}_1) \cdot \pairing(\bar{\tid}^{1/\delta}, \widetilde{u}_2) =
\pairing(\bar{\gen}^{1/\delta}, \ggen)= \pairing(\gen, \ggen) \ \wedge \\
  %
    \pairing(\bar{\signaturetwo}^{1/\delta}, \ggen) =
\pairing(\bar{\tid}^{1/\delta},
\signatureone')\cdot\prod_{i\in\ISSUERS}\pairing(H_i^{1/\delta}, \ipk{i}).
  %
  \end{gather*}
%
  So the following hold:
  \begin{itemize}
    \item $(\signatureone', \bar{\signaturetwo}^{1/\delta})$ is a valid
tag-based aggregate signature on the preimage of $\bar{\tid}^{1/\delta}$ and
the attribute values $\atv{i}$;
    \item $(\widetilde{r}', \widetilde{s}', \bar{z}^{1/\delta})$ is a valid
$\AGHO$ signature on $(\pk'^{1/\delta}, \bar{\tid}^{1/\delta})$.
  \end{itemize}

  To extract the preimage of $\bar{\tid}^{1/\delta}$, we use the following:

  \begin{enumerate}

    \item Random oracle $\hashtwo$, which is initialised with two maps $L_1$ and
$L_2$, and which for each query $\Tag$, checks if $L_1[\Tag]$ is empty. If so,
it  selects a random group element $\tid \in \group{1}$, and sets $L_1[\Tag]
\leftarrow \tid$ and $L_2[\tid] \leftarrow \Tag$.  Finally, returns $L_1[\Tag]$.

    \item The fact that the registration authority only signs well-formed pairs
defined as $(\pk, \hashtwo(\Tag))$.

  \end{enumerate}

  By existential unforgeability of $\AGHO$, the registration authority signed
the pair $(\pk'^{1/\delta}, \bar{\tid}^{1/\delta})$. It signs only well-formed
pairs, so $\bar{\tid}^{1/\delta}$ is the output of a query to $\hashtwo$, and
the tag is its entry in $L_2$. If the adversary also sees simulated
presentations, the reduction to unforgeability of $\AGHO$ has no $\rask$. It
answers each simulated presentation with a signing query on the random pair
$(\pk'^{1/\theta}, \bar{\tid}^{1/\theta})$, in the notation of the simulation
below, and sets $\bar{z} = z^\theta$ for the returned $z$. An extracted pair
equal to one of these gives $\sk$ with $\gen^\sk = \pk'^{1/\theta}$, a discrete
logarithm of a uniform element of $\group{1}$. The extracted values then give
the following:
  \begin{itemize}
    \item the tag $\Tag$, with $\tid = \bar{\tid}^{1/\delta}$;
    \item the secret key $\sk$, with $\pk = \gen^\sk$;
    \item the $\AGHO$ signature $\SPSsig = (\widetilde{r}', \widetilde{s}',
\bar{z}^{1/\delta})$ on $(\pk, \tid)$;
    \item the tag-based aggregate signature $\signature = (\signatureone',
\bar{\signaturetwo}^{1/\delta})$ on $\tid$ and the attribute values $\atv{i}$.
  \end{itemize}
Together these are a witness $\witness_2 = (\sk, \pk, \tid, \signature,
\SPSsig)$ for $\rel_2$, and with $\Tag$ a witness $(\Tag, \sk, \pk, \SPSsig,
\signature)$ for $\rel_1$, with $\DS$ set to $\AGHO$ on $(\pk, \hashtwo(\Tag))$.

  \paragraph{Simulation.} If DDH holds in $\group{1}$, a simulator holding
$\rask = (\mu_1, \mu_2, \tau)$ can simulate a valid presentation $\algoproof$
without a witness.

  \begin{enumerate}

    \item Randomly select $\theta, \omega, \varphi, \nu, \{\epsilon_i\}_{i\in
\ISSUERS},  \iota \in \Zp^*$, and set the following:

%
      \begin{equation*}
%
	\bar{\gen} = \gen^\theta, \quad \bar{v} = v^\theta, \quad
\bar{\signaturetwo} = \gen^\varphi, \quad \bar{\tid} = \gen^\nu, \quad H_i =
\gen^{\epsilon_i}, \quad \pk' = \gen^\iota.
%
      \end{equation*}
%
      In an honest presentation, $\bar{\gen} = \gen^\delta$ is uniform, and
$\bar{v}, \bar{\tid}, \pk'$ and the $H_i$ are the $\delta$-th powers of $v,
\tid, \pk$ and the $\hashone(\atv{i})$. Here $\bar{v}$ is computed in the same
way, with $\theta$ in place of $\delta$. A hybrid argument that replaces
$\bar{\tid}$, $\pk'$ and the $H_i$ by uniform elements one at a time, each step
a DDH instance in $\group{1}$, shows that they are indistinguishable from the
uniform elements above.

    \item We define the signature elements as follows:
%
      \begin{gather*}
%
	\widetilde{r}' = \ggen^\omega, \quad \widetilde{s}' =
\widetilde{r}'^{\,\tau}, \quad \bar{z} = (\bar{\gen}\,\pk'^{-\mu_1}
\bar{\tid}^{-\mu_2})^{1/\omega}, \\ \signatureone' =
(\ggen^{\varphi}\textstyle\prod_{i \in \ISSUERS}\ipk{i}^{-\epsilon_i})^{1/\nu}.
%
      \end{gather*}
%
      In an honest presentation, $(\widetilde{r}', \widetilde{s}', \bar{z})$
is a randomised $\AGHO$ signature on $(\pk, \tid)$ with $\bar{z}$ raised to
$\delta$. This is a fresh $\AGHO$ signature on $(\pk', \bar{\tid})$ with
$\bar{\gen}$ in place of $\gen$, which the simulator computes with $\rask$.
$\signatureone'$ is uniform after randomisation, and $\bar{\signaturetwo}$ is
fixed by the verification equation below. Here $\bar{\signaturetwo}$ is uniform
and $\signatureone'$ is fixed by the same equation, which gives the same
distribution.
      Now we have by construction that $(\bar{\gen}, \bar{v}, \widetilde{r}',
\widetilde{s}', \bar{z}, \signatureone', \bar{\signaturetwo}, \bar{\tid}, \pk',
\{\atv{i}, H_i\}_{i\in \ISSUERS})$ satisfies the following:
    %
      \begin{gather*}
    %
	\pairing(\bar{\signaturetwo}, \ggen) =
\prod_{i\in\ISSUERS}\pairing(H_i,\ipk{i}) \cdot
\pairing(\bar{\tid},\signatureone'), \\
      %
	\pairing(\bar{\gen},\ggen) = \pairing(\bar{z},\widetilde{r}')\cdot
\pairing(\pk',\widetilde{u}_1) \cdot \pairing(\bar{\tid},\widetilde{u}_2), \\
      %
	\pairing(\bar{v},\widetilde{r}') = \pairing(\bar{\gen},\widetilde{s}').
    %
      \end{gather*}
    %

    \item Pick $\pi_1$, $\pi_2$ and $c$ randomly in $\Zp$, and set the
following:
%
      \begin{equation*}
%
	f = \gen^{\pi_1}\cdot\bar{\gen}^{-c}, \quad h =
v^{\pi_1}\cdot\bar{v}^{-c}, \quad y_i = \hashone(\atv{i})^{\pi_1}\cdot H_i^{-c},
\quad w = \bar{\gen}^{\pi_2}\cdot\pk'^{-c}.
%
      \end{equation*}

        \item Program random oracle $\hash$ so that $\hash(\gen, v, \inst_3, f, h,
\{y_i\}_{i \in \ISSUERS}, w, \msg) = c$, where $\inst_3$ is the simulated
statement above.

    \item Finally, set $\eta$ to tuple $(f, h, \{y_i\}_{i \in \ISSUERS}, w,
\pi_1, \pi_2)$.

  \end{enumerate}

  So under the DDH assumption in $\group{1}$, the simulated and honestly
generated presentations are computationally indistinguishable. This proves
\Cref{lem:agg-creds-sigma}.

\end{proof}

\subsection{Proof of \Cref{thm:agg-creds-thm-2}: security of the instantiation}\label{sec:agg-creds-proof-3}

The instantiation is an instance of the construction that sets $\DS$ to the
$\AGHO$ SPS scheme, $\ATS$ to our tag-based aggregate signature scheme, and the
proof system to the presentation proof of $\AACpresent$ and $\AACverify$.

\begin{proof}

  \begin{description}

        \item[$\DS$ is existentially unforgeable.] The $\AGHO$ SPS scheme is existentially
unforgeable under chosen-message attack in the GGM~\cite{C:AGHO11}.

    \item[$\ATS$ is existentially unforgeable under tag-based aggregation.] This
is \Cref{thm:agg-creds-thm-1}, under the BCDH assumption in the random oracle
model.

    \item[The proof system is zero knowledge and simulation extractable.] A
proof for $\inst_1$ is $\algoproof = (\eta, \inst_3)$, verified by
$\AACverify$. The Sigma protocol for $\rel_3$ of
\eqref{eq:agg-creds-proof-rel} has special soundness and honest-verifier zero
knowledge. Its responses are unique, and so quasi-unique, since each response
is fixed by one verification equation. The challenge hashes the whole statement
together with $\msg$. The Fiat-Shamir transform then makes the proof for
$\rel_3$ simulation extractable with $\msg$ as its label in the random oracle
model~\cite{INDOCRYPT:FKMV12}. \Cref{lem:agg-creds-sigma} turns this into
extraction of a witness for $\rel_1$, except with the probability of forging
$\AGHO$ or computing a discrete logarithm in $\group{1}$, both negligible, and
simulation of presentations under the DDH assumption in $\group{1}$. The
simulator programs the random oracle to simulate honest users' presentations.
It needs $\rask$, which the simulator of the construction's proof holds, since
it plays the honest registration authority. The
simulator of the construction's proof runs straight-line and the
reductions for $a_1$, $a_2$ and $a_3$ extract, each once, at the first
presentation where $E$ occurs. Each
extraction loses a polynomial factor in the number of hash
queries~\cite{JC:PoiSte00}.

    \item[$f$ is one way, and the proofs of knowledge of secret keys are
zero-knowledge proofs of knowledge.] Here $f(\sk) = \gen^\sk$, which is one way
under the DL assumption in $\group{1}$.
Users' and issuers' proofs of knowledge are Sigma protocols for a discrete
logarithm with the Fiat-Shamir transform, which are zero knowledge and proofs of knowledge in the random
oracle model~\cite{JC:PoiSte00}.

  \end{description}

  The instantiated protocol securely realises ideal functionality
$\idfu{\AAC}$ in the GGM and the random oracle model.

\end{proof}

